\chapter{Počítačová paměť}

K uchovávání většího objemu dat, než může dovolit počet procesoru, slouží \textbf{počítačová paměť.}
Toto zařízení skládá se z \textbf{paměťových buněk}, každá z kterých uchovává jednotku informace.
Každá buňka má pevně přiřazenou adresu, pomocí které procesor může k této buňce přistupovat
(adresovat). Paměťová buňka je minimální adresovací jednotka.

Rozdělení pamětí podle :

\begin{itemize}
    \item \textbf{Random access memory} (\textit{RAM}) - paměť, podporující zápis i čtení.
        \begin{itemize}
            \item Static RAM (SRAM) - každá paměťová buňka je sestavena z D flip-flopů.
            \item Dynamic RAM (DRAM) - každá paměťová buňka je sestavena z tranzistoru (zpravidla
                MOSFET) a kondenzátoru. Hlavní nevýhodou je, že nahromaděný náboj na kondenzátoru má
            tendenci unikat, kvůli čemu táto paměť musí se periodicky obnovovat (\textit{refresh}).
            \item Synchronous DRAM (SDRAM) - hybrid statické a dynamické RAM.
        \end{itemize}
    \item \textbf{Read-only memory} (\textit{ROM}) - paměť, podporující jenom čtení. Obsah je
        předem zapsán výrobcem pamětí.
\end{itemize}

Obvykle rozdělení RAM a ROM podle technologie nemá význam pro emulaci, ale pro větší přesnost
emulátor může např. emulovat  unik elektrických nábojů v DRAM pamětí.

\section{Reprezentace dat v počítači}

Většina počítačových systémů využívá \emph{digitální veličiny} k reprezentaci vnitřního stavu a dalších
dat. Digitální veličiny jsou hodnoty, které nabývají konečného počtu diskrétních stavů. Tento přístup
vychází z fyzikálních vlastností elektrických obvodů. Přestože elektrické obvody pracují s
analogovými veličinami, jako jsou napětí a proud, lze tyto veličiny rozdělit do určitých intervalů,
přičemž každý interval odpovídá jednomu diskrétnímu stavu. S rostoucím počtem rozlišovaných stavů
se však zmenšuje rozdíl mezi sousedními hodnotami, čímž se systém stává méně spolehlivým vlivem
elektrického šumu. Zároveň je nutné používat přesnější a složitější elektrické obvody, což zvyšuje
jejich cenu a náročnost návrhu. Z těchto důvodů se pro reprezentaci digitálních hodnot používá
dvojková číselná soustava, ve které je nutné rozlišit pouze dva stavy — 0 a 1.

Jedna jednotka takové informace, která nabývá stavů 0 a 1, se nazývá \textbf{bit} (z anglického
\textit{binary digit}). Další základní jednotkou je \textbf{bajt} neboli \textbf{byte} o velikosti 8
bitů, a nabývá hodnot 0 až 255 v desítkové soustavě. Kombinace dvou bytů (16 bitů) tvoří \textbf{slovo}
(\textit{word}). Kombinace 4 bytů (32 bitů) tvoří dvouslovo (\textit{double word} --
\textit{dword}).

V tabulce \ref{tab:data_units} je uveden přehled jednotek odvozených od základní jednotky byte.
Historicky se používaly předpony soustavy SI, které vyjadřují mocniny deseti. Pro zachování vazby na
mocniny dvou byly později zavedeny binární předpony, jako například kibi nebo mebi.

\begin{table}
    \centering
    \begin{tabular}{llp{5cm}l}
        \toprule
        \textbf{Jednotka} & \textbf{Značka} & \textbf{B} & Mocnina $2^n$ \\
        \midrule
        Byte & B & 8 & $2^3$ \\
        Kilobyte & kB & \num{1000} & -- \\
        Kibibyte & KiB & \num{1024} & $2^{10}$ \\
        Megabyte & MB & \num{1000000} & -- \\
        Mebibyte & MiB & \num{1048576} & $2^{20}$ \\
        Gigabyte & GB & $10^9$ & -- \\
        Gibibyte & GiB & ~$\num{1.074}\times10^9$ & $2^{30}$\\
        Terabyte & TB & $10^{12}$ & -- \\
        Tebibyte & TiB & ~$\num{1.1}\times10^{12}$ & $2^{40}$ \\
        \bottomrule
    \end{tabular}
    \caption{Přehled jednotek digitálních dat.}
    \label{tab:data_units}
\end{table}

Historicky, u různých komerčních počítačů jedna paměťová buňka mohla obsahovat různý počet bitů, ale
později se ustálilo použití jednoho bytu (8 bitů) jako minimální adresovací jednotky v pamětí.

\section{Adresní prostor}

\textbf{Adresní prostor} je rozsah diskrétních paměťových adres, přičemž každá jednotlivá adresa
muže odpovídat zařízení, diskovému sektoru, registru, fyzické paměti, nebo jiné entitě. Díky
adresnímu prostoru, paměť slouží pro CPU nejenom jako úložiště dat, ale také jako \textbf{mapa
okolního světa}.

\begin{figure}[htbp]
    \centering
    \begin{subfigure}{0.45\textwidth}
        \includesvg[width=0.9\textwidth]{images/flat_memory.svg}
        \caption{Lineární adresní prostor.}
        \label{fig:flat_memory}
    \end{subfigure}
    \quad
    \begin{subfigure}{0.45\textwidth}
        \includesvg[width=1.1\textwidth]{images/segmented_memory.svg}
        \caption{Segmentovaný adresní prostor.}
        \label{fig:segmented_memory}
    \end{subfigure}
    \caption{Reprezentace adresního prostoru.}
    \label{fig:address_space}
\end{figure}

Existují dva rozšířené modely adresního prostoru:

\begin{itemize}
    \item \textbf{Lineární adresní prostor (ploché adresní rozložení)} -- metoda adresování pamětí,
        kde je procesoru přímo přístupný celý rozsah pamětí jako jedno souvislé jednorozměrné pole.
        Každá buňka má unikátní index od $0$ do $2^n - 1$, kde $n$ je bitová šířka sběrnice (obr.
        \ref{fig:flat_memory}).
    \item \textbf{Segmentování pamětí} -- model, který rozděluje adresní prostor na logické celky zvané
        \textbf{segmenty}. Adresa se pak neskládá z jednoho čísla, ale z dvojice: \textit{selektor
        segmentu} a \textit{offset} (posun) v rámci tohoto segmentu. Tato metoda byla klíčová pro
        překonání limitů 16-bitových procesorů (např. Intel 8086), kde umožnila adresovat až 1~MB
        paměti pomocí dvou 16-bitových hodnot (obr. \ref{fig:segmented_memory}).
\end{itemize}

\section{Bank switching}

U mnoha historických 8-bitových systémů (např. NES) se setkáváme s problémem, kdy je
adresní rozsah procesoru (typicky 64~KiB) nedostatečný pro rostoucí nároky softwaru. Řešením je
technika \textbf{přepínání bank (bank switching)}.

Tento mechanismus umožňuje procesoru přistupovat k mnohem větší fyzické paměti tím, že v adresní
mapě vytvoří „okno“ (banku), do kterého lze dynamicky mapovat různé části externí paměti.

\section{Mapování}

\begin{figure}[H]
    \centering
    \includesvg[width=\textwidth]{images/memory_map.svg}
    \caption{Mapovaní pamětí: RAM, grafická jednotka, zvuková jednotka a vstup.}
    \label{fig:memory_map}
\end{figure}

Přiřazení rozsahu adresního prostoru nějakému zařízení se nazývá \textbf{mapování}, a je
zprostředkováno primárně sběrnici (viz kapitola \ref{ch:tp-bus}). Přičemž z pohledu CPU neexistuje
rozdíl mezi přístupem k pamětí nebo k zařízení - existují pouze paměťové adresy, a pro komunikaci
s externími zařízeními se používají stejné instrukce jako pro běžný přístup k pamětí. I paměť je
externí zařízení, přípojené k sběrnici (obr. \ref{fig:memory_map}).

Jakýkoli Přístup k pamětí může mít \textbf{vedlejší účinky} (\textit{side effects}) na CPU nebo
okolí komponenty (včetně přistupovaného). Např. v systému NES čtení stavového registru grafické
jednotky dojde k vynulovaní příznaku sudosti aktuálního zápisu do video pamětí.

\section{Endianita (pořadí bajtů)}

\textbf{Endianita} definuje způsob, jakým jsou v paměti uloženy vícebytové datové typy (např.
16-bitová nebo 32-bitová celá čísla). Rozlišujeme dva základní typy:

\begin{itemize}
    \item \textbf{Little-endian}: Nejméně významný byte (LSB) je uložen na nejnižší adrese. Toto
        schéma využívá většina moderních procesorů (x86, ARM).
    \item \textbf{Big-endian}: Nejvýznamnější byte (MSB) je uložen na nejnižší adrese. Tento způsob
        je přirozenější pro lidské čtení hexadecimálních výpisů a využívají jej například procesory
        řady Motorola 68000.
\end{itemize}
